\chapter{中值定理}
\label{chap:mean-value-theorems}

导数是在一点取得的局部信息。若要由导数判断一段区间上的增减、估计两个函数值之差，或者比较两个趋于零或无穷的函数，就需要把局部信息同区间两端的增量联系起来。中值定理完成的正是这一步：某个中间点的导数再现或控制整段区间的平均变化率。

本章中的函数均为实值函数，除非另有说明。中值定理只要求闭区间上的连续性和开区间内的可导性；端点处不必可导。这个条件分工在后面的反例和 L'Hôpital 法则中都会起作用。

\section{Fermat、Rolle 与 Lagrange 定理}
\label{sec:fermat-rolle-lagrange}

\subsection{内点极值与 Fermat 定理}

设 \(f:D\to\R\)，\(a\in D\)。若存在 \(\delta>0\)，使
\[
 f(x)\le f(a)
 \qquad\bigl(x\in D,\ \abs{x-a}<\delta\bigr),
\]
则称 \(a\) 为 \(f\) 的局部极大点；局部极小点对称定义。这里的“局部”只比较定义域中充分靠近 \(a\) 的点。若 \(a\) 是区间的内点，则可以从左右两侧考察差商。

\begin{theorem}[Fermat 定理]
\label{thm:fermat-stationary}
设 \(a\) 是区间 \(I\) 的内点，函数 \(f:I\to\R\) 在 \(a\) 可导。若 \(a\) 是 \(f\) 的局部极大点或局部极小点，则
\[
 f'(a)=0.
\]
\end{theorem}

\begin{proof}
只证局部极大情形。存在 \(\delta>0\)，使 \(x\in I\)、\(\abs{x-a}<\delta\) 时 \(f(x)\le f(a)\)。对充分小的 \(h>0\)，
\[
 \frac{f(a+h)-f(a)}{h}\le0;
\]
对充分小的 \(h<0\)，分子仍不大于零而分母为负，故
\[
 \frac{f(a+h)-f(a)}{h}\ge0.
\]
因 \(f'(a)\) 存在，左右差商必须趋于同一个数。右极限不大于零，左极限不小于零，所以共同极限只能为零。局部极小情形把不等号反向即可。
\end{proof}

满足 \(f'(a)=0\) 的点称为驻点。Fermat 定理给出内点可导极值的必要条件，却不是充分条件；例如 \(f(x)=x^3\) 在 \(0\) 处导数为零，但 \(0\) 不是局部极值点。它也没有处理端点极值或尖点极值，这两种情形将在本章第四节中逐一核对。

\subsection{Rolle 定理}

\begin{theorem}[Rolle 定理]
\label{thm:rolle}
设 \(a<b\)，函数 \(f:[a,b]\to\R\) 在 \([a,b]\) 上连续，在 \((a,b)\) 内可导，并且
\[
 f(a)=f(b).
\]
则存在 \(c\in(a,b)\)，使 \(f'(c)=0\)。
\end{theorem}

\begin{proof}
闭区间连续函数取得最大值和最小值。若二者相等，则 \(f\) 为常值，任取 \(c\in(a,b)\) 都有 \(f'(c)=0\)。

若最大值与最小值不同，由 \(f(a)=f(b)\)，二者不可能都只在两个端点取得：若最大值恰为端点公共值且函数不恒定，则最小值在某个内点取得；若最小值恰为端点公共值，则最大值在某个内点取得；若端点公共值介于最大值与最小值之间，两个极值都在内点取得。因此总有某个 \(c\in(a,b)\) 是局部极值点。Fermat 定理给出 \(f'(c)=0\)。
\end{proof}

证明中闭区间的作用不是装饰性的：它使连续函数取得极值。端点等值再迫使至少一个极值落入可以使用 Fermat 定理的区间内部。

\subsection{Lagrange 中值定理}

\begin{theorem}[Lagrange 中值定理]
\label{thm:lagrange-mvt}
设 \(a<b\)，函数 \(f:[a,b]\to\R\) 在 \([a,b]\) 上连续，在 \((a,b)\) 内可导。则存在 \(c\in(a,b)\)，使
\[
 f(b)-f(a)=f'(c)(b-a),
\]
即
\[
 f'(c)=\frac{f(b)-f(a)}{b-a}.
\]
\end{theorem}

\begin{proof}
令
\[
 \ell(x)=f(a)+\frac{f(b)-f(a)}{b-a}(x-a)
\]
为连接两端点的割线函数，并设 \(F=f-\ell\)。函数 \(F\) 在 \([a,b]\) 上连续、在 \((a,b)\) 内可导，而且
\[
 F(a)=F(b)=0.
\]
由 Rolle 定理，存在 \(c\in(a,b)\) 使 \(F'(c)=0\)。因此
\[
 f'(c)-\frac{f(b)-f(a)}{b-a}=0,
\]
整理即得结论。
\end{proof}

Lagrange 定理把端点增量写成一个中间点导数与区间长度的乘积。这里的 \(c\) 依赖于函数和两个端点，通常不能预先指定，也不必唯一。把 \(a\) 固定、令 \(b=a+h\)，可写成有限增量形式
\[
 f(a+h)-f(a)=f'(a+\theta h)h
\]
其中某个 \(\theta\in(0,1)\) 依赖于 \(h\)。这个公式是精确等式，但它并没有说明 \(\theta\) 随 \(h\) 怎样变化。

\begin{example}[指数函数的基本估计]
对 \(x>0\)，把 Lagrange 定理用于 \(e^t\) 在 \([0,x]\) 上的限制，得到某个 \(c\in(0,x)\) 使
\[
 e^x-1=e^c x.
\]
于是
\[
 x\le e^x-1\le xe^x.
\]
对负 \(x\) 可在 \([x,0]\) 上作同样分析。与只写出导数公式相比，中值定理保留了区间上的定量控制。
\end{example}

\section{Cauchy 中值定理}
\label{sec:cauchy-mean-value}

Lagrange 定理比较函数值增量与自变量增量。若把自变量 \(x\) 换成另一个函数 \(g(x)\)，就得到两个函数增量之间的比较。这一形式是 L'Hôpital 法则的直接基础。

\begin{theorem}[Cauchy 中值定理]
\label{thm:cauchy-mvt}
设 \(a<b\)，函数 \(f,g:[a,b]\to\R\) 都在 \([a,b]\) 上连续、在 \((a,b)\) 内可导。则存在 \(c\in(a,b)\)，使
\[
 \bigl(f(b)-f(a)\bigr)g'(c)
 =
 \bigl(g(b)-g(a)\bigr)f'(c).
\]
\end{theorem}

\begin{proof}
定义
\[
 H(x)=\bigl(f(b)-f(a)\bigr)\bigl(g(x)-g(a)\bigr)
      -\bigl(g(b)-g(a)\bigr)\bigl(f(x)-f(a)\bigr).
\]
它在闭区间上连续、在开区间内可导，并且直接计算得
\[
 H(a)=H(b)=0.
\]
Rolle 定理给出 \(c\in(a,b)\) 使 \(H'(c)=0\)。展开这一等式便得到
\[
 \bigl(f(b)-f(a)\bigr)g'(c)
 -
 \bigl(g(b)-g(a)\bigr)f'(c)=0.
\]
\end{proof}

交叉相乘的陈述始终有意义，不要求 \(g(b)\ne g(a)\)，也不要求 \(g'(c)\ne0\)。若要写成商，必须另加非零条件。

\begin{corollary}[增量比形式]
\label{cor:cauchy-mvt-ratio}
在 Cauchy 中值定理的条件下，若
\[
 g'(x)\ne0\qquad(x\in(a,b)),
\]
则 \(g(b)\ne g(a)\)，并且存在 \(c\in(a,b)\) 使
\[
 \frac{f(b)-f(a)}{g(b)-g(a)}
 =\frac{f'(c)}{g'(c)}.
\]
\end{corollary}

\begin{proof}
若 \(g(a)=g(b)\)，Rolle 定理会给出某个 \(\xi\in(a,b)\) 使 \(g'(\xi)=0\)，与假设矛盾。因此端点增量非零。再在 Cauchy 中值定理的等式两端除以
\((g(b)-g(a))g'(c)\) 即得商式。
\end{proof}

\begin{corollary}[导数比控制增量比]
\label{cor:derivative-ratio-controls-increments}
设上述非零条件成立。若对所有 \(x\in(a,b)\) 有
\[
 m\le\frac{f'(x)}{g'(x)}\le M,
\]
则
\[
 m\le
 \frac{f(b)-f(a)}{g(b)-g(a)}
 \le M.
\]
\end{corollary}

\begin{proof}
增量比等于某个中间点 \(c\) 处的导数比，代入给定上下界即可。这个结论不要求 \(g'\) 为正；分母可能为负的影响已经包含在两个比值中。
\end{proof}

取 \(g(x)=x\) 时，Cauchy 中值定理退化为 Lagrange 中值定理。取一般 \(g\) 则允许用 \(g\) 的变化作为尺度，这正是比较 \(f/g\) 型极限所需的形式。

\section{单调性、常值性与导数估计}
\label{sec:mvt-consequences}

\subsection{导数符号与单调性}

\begin{theorem}[导数符号判据]
\label{thm:derivative-sign-monotonicity}
设 \(I\) 是区间，\(f:I\to\R\) 在 \(I\) 的内点可导，并在属于 \(I\) 的端点连续。
\begin{enumerate}
 \item 若 \(f'(x)\ge0\) 对所有内点成立，则 \(f\) 在 \(I\) 上单调不减；
 \item 若 \(f'(x)\le0\) 对所有内点成立，则 \(f\) 在 \(I\) 上单调不增；
 \item 若 \(f'(x)>0\) 对所有内点成立，则 \(f\) 严格递增；若 \(f'(x)<0\)，则 \(f\) 严格递减。
\end{enumerate}
反过来，若 \(f\) 在 \(I\) 上单调不减且在内点 \(x\) 可导，则 \(f'(x)\ge0\)；单调不增时 \(f'(x)\le0\)。
\end{theorem}

\begin{proof}
任取 \(x<y\) 于 \(I\)。函数在 \([x,y]\) 上连续、在 \((x,y)\) 内可导，故存在 \(c\in(x,y)\) 使
\[
 f(y)-f(x)=f'(c)(y-x).
\]
由于 \(y-x>0\)，前三项结论由 \(f'(c)\) 的符号立即得到。

反过来，设 \(f\) 单调不减，且 \(x\) 为内点。对充分小的 \(h\ne0\)，当 \(h>0\) 时分子 \(f(x+h)-f(x)\ge0\)；当 \(h<0\) 时分子不大于零，除以负数后差商仍不小于零。差商极限存在时必有 \(f'(x)\ge0\)。单调不增情形对称。
\end{proof}

严格递增函数的导数可以在某些点等于零，例如 \(x^3\) 在 \(0\) 处。因而“严格递增”不能反推 \(f'>0\) 处处成立；正确的反向结论仍只是 \(f'\ge0\)。

在已经知道 \(f'\ge0\) 时，严格递增仍有一个准确判据，它比处处 \(f'>0\) 弱。

\begin{proposition}[非负导数与严格递增]
\label{prop:nonnegative-derivative-strictness}
在 \theoremref{thm:derivative-sign-monotonicity} 的条件下，设 \(f'\ge0\) 于 \(I\) 的全部内点。则 \(f\) 严格递增，当且仅当 \(f'\) 不在 \(I\) 的任何非退化子区间内部恒为零。
\end{proposition}

\begin{proof}
若某个非退化子区间上 \(f'\equiv0\)，零导数判据说明 \(f\) 在该子区间为常值，故不可能严格递增。

反过来，导数非负已经保证 \(f\) 单调不减。若存在 \(x<y\) 使 \(f(x)=f(y)\)，则对每个 \(z\in[x,y]\) 都有
\[
 f(x)\le f(z)\le f(y)=f(x).
\]
所以 \(f\) 在 \([x,y]\) 上恒定，进而 \(f'=0\) 于 \((x,y)\)，与假设矛盾。因此 \(x<y\) 总有 \(f(x)<f(y)\)。
\end{proof}

\begin{corollary}[零导数判据]
\label{cor:zero-derivative-constant}
设 \(I\) 是区间，\(f:I\to\R\) 满足上一结论的连续与可导条件。则
\[
 f'(x)=0\quad(x\in\operatorname{int}I)
\]
当且仅当 \(f\) 在 \(I\) 上为常值。
\end{corollary}

\begin{proof}
若导数恒为零，则它既非负又非正，故 \(f\) 同时单调不减和单调不增，只能为常值。常值函数的导数显然处处为零。
\end{proof}

\subsection{增量估计}

\begin{theorem}[导数上下界与增量]
\label{thm:derivative-bounds-increments}
设 \(f:[a,b]\to\R\) 连续并在 \((a,b)\) 内可导。若
\[
 m\le f'(x)\le M\qquad(x\in(a,b)),
\]
则对任意 \(a\le x<y\le b\)，
\[
 m(y-x)\le f(y)-f(x)\le M(y-x).
\]
\end{theorem}

\begin{proof}
把 Lagrange 中值定理用于 \([x,y]\)，得到某个 \(c\in(x,y)\) 使
\[
 f(y)-f(x)=f'(c)(y-x).
\]
因 \(y-x>0\)，将 \(m\le f'(c)\le M\) 同乘以 \(y-x\) 即得。
\end{proof}

\begin{corollary}[有界导数与 Lipschitz 估计]
\label{cor:bounded-derivative-lipschitz}
若
\[
 \abs{f'(x)}\le K\qquad(x\in(a,b)),
\]
则
\[
 \abs{f(y)-f(x)}\le K\abs{y-x}
\qquad(x,y\in[a,b]).
\]
因此 \(f\) 是 \(K\)-Lipschitz 连续。反之，若 \(f\) 在区间 \(I\) 上是 \(K\)-Lipschitz 连续，并在内点 \(x\) 可导，则
\[
 \abs{f'(x)}\le K.
\]
\end{corollary}

\begin{proof}
前一结论取 \(m=-K\)、\(M=K\)。反向中，对充分小的非零 \(h\) 有
\[
 \abs{\frac{f(x+h)-f(x)}{h}}\le K.
\]
令 \(h\to0\)，得到 \(\abs{f'(x)}\le K\)。
\end{proof}

这个等价关系说明，在一元区间上，“导数处处受统一界控制”和“函数具有全局 Lipschitz 增量估计”是同一项定量事实的微分与差分表述。若只知道每一点的导数有限，却没有统一上界，就不能推出整个区间上的同一 Lipschitz 常数。

这里不需要区间有界或闭。若 \(I\) 是任意区间，\(f\) 在内点可导、在属于 \(I\) 的端点连续，并满足 \(|f'|\le K\)，则任取 \(x<y\) 于 \(I\)，在闭区间 \([x,y]\) 上应用中值定理即可得到同一 Lipschitz 估计。反向结论也只使用点附近的差商。因此在这样的区间上，导数的统一界与 Lipschitz 常数可以逐点相互恢复。

\begin{corollary}[导数比较]
\label{cor:derivative-comparison}
若 \(f,g:[a,b]\to\R\) 连续、在 \((a,b)\) 内可导，且 \(f'\le g'\)，则对任意 \(a\le x<y\le b\)，
\[
 f(y)-f(x)\le g(y)-g(x).
\]
\end{corollary}

\begin{proof}
函数 \(h=f-g\) 满足 \(h'\le0\)，故单调不增。
\end{proof}

\begin{proposition}[下导数界与逆向距离估计]
\label{prop:lower-derivative-inverse-lipschitz}
设 \(f:[a,b]\to\R\) 连续、在 \((a,b)\) 内可导，并且或者
\[
 f'(x)\ge m>0,
\]
或者 \(f'(x)\le-m<0\)。则 \(f\) 严格单调，且对所有 \(x,y\in[a,b]\)，
\[
 \abs{f(x)-f(y)}\ge m\abs{x-y}.
\]
因而逆函数 \(f^{-1}:f([a,b])\to[a,b]\) 是 \(1/m\)-Lipschitz 连续。
\end{proposition}

\begin{proof}
严格单调性来自导数符号判据。对 \(x<y\)，Lagrange 定理给出
\[
 \abs{f(y)-f(x)}=\abs{f'(c)}(y-x)\ge m(y-x).
\]
令 \(u=f(x)\)、\(v=f(y)\)，改写为
\[
 \abs{f^{-1}(u)-f^{-1}(v)}
 \le\frac1m\abs{u-v}.
\]
\end{proof}

\section{各项假设的作用和反例}
\label{sec:mvt-hypotheses-counterexamples}

中值定理的假设分属三个位置：闭区间连续性把端点同内部连接起来，开区间可导性使 Fermat 定理能够作用于内部极值，区间性则保证两端之间的所有点都属于定义域。删去其中任何一项，都可能失去中间点。

\subsection{Fermat 定理中的三个条件}

\begin{counterexample}[缺少内点条件]
函数 \(f(x)=x\) 在 \([0,1]\) 的左端点 \(0\) 取得最小值，但右导数为 \(1\)。Fermat 定理只对区间内点作结论；端点极值只给出单侧差商的符号。
\end{counterexample}

\begin{counterexample}[缺少可导性]
函数 \(f(x)=\abs x\) 在 \(0\) 处取得严格局部最小值，但左右导数分别为 \(-1\) 和 \(1\)，导数不存在。尖点极值不能由 \(f'(a)=0\) 描述。
\end{counterexample}

\begin{counterexample}[驻点不一定是极值点]
函数 \(f(x)=x^3\) 满足 \(f'(0)=0\)，但它在 \(0\) 的任意邻域内同时取正值和负值。因此 Fermat 条件不可反向使用。
\end{counterexample}

\subsection{Rolle 与 Lagrange 定理中的条件}

\begin{counterexample}[闭区间连续性不可删除]
在 \([0,1]\) 上定义
\[
 f(0)=1,\qquad f(x)=x\quad(0<x\le1).
\]
函数在 \((0,1)\) 内可导且 \(f(0)=f(1)=1\)，但 \(f'(x)=1\) 对所有 \(x\in(0,1)\) 成立。由于函数在 \(0\) 处不连续，Rolle 定理的结论失败；同一例也使端点割线斜率为零而内部导数从不为零，故 Lagrange 定理同样失败。
\end{counterexample}

\begin{counterexample}[内部可导性不可删除]
函数 \(f(x)=\abs x\) 在 \([-1,1]\) 上连续且端点值相等，但内部没有导数为零的点：在负半轴导数为 \(-1\)，在正半轴导数为 \(1\)，而 \(0\) 处不可导。因此 Rolle 定理不能只假设连续。它的端点割线斜率为零，也说明 Lagrange 定理需要整个开区间内可导。
\end{counterexample}

\begin{counterexample}[端点等值只属于 Rolle 定理]
函数 \(f(x)=x\) 在 \([0,1]\) 上连续、在内部可导，却没有导数为零的点。这里 \(f(0)\ne f(1)\)，所以不能使用 Rolle 定理；Lagrange 定理仍成立，因为内部导数恒等于割线斜率 \(1\)。
\end{counterexample}

\begin{counterexample}[定义域必须连接两端]
令
\[
 D=[-1,0]\cup[1,2],
\qquad
 f(x)=
 \begin{cases}
 x+1,&x\in[-1,0],\\
 x-2,&x\in[1,2].
 \end{cases}
\]
函数在相对定义域 \(D\) 上连续，在两个分支内部的导数都为 \(1\)，并且 \(f(-1)=f(2)=0\)。若把 Rolle 定理错误推广到不连通定义域，就会预言某处导数为零。原定理要求函数定义在整个 \([a,b]\) 上，使端点之间的极值与割线论证能够进行。
\end{counterexample}

端点可导性则不是中值定理的条件。例如 \(f(x)=\sqrt x\) 在 \([0,1]\) 上连续、在 \((0,1)\) 内可导，虽然在 \(0\) 处没有有限右导数，Lagrange 定理仍然适用。把端点可导加入假设会得到一个正确但不必要地偏强的版本。

Cauchy 中值定理取 \(g(x)=x\) 就退化为 Lagrange 定理，所以上述闭区间连续性、内部可导性和区间性反例也说明 Cauchy 定理的相应条件不能随意删除。另一方面，交叉相乘形式本身不要求 \(g'\ne0\)；这一非零条件只在把结论改写成导数比与增量比时加入。

\begin{remark}[严格符号与非严格单调性]
由 \(f'\ge0\) 只能推出单调不减，不能推出严格递增；常值函数就是反例。由 \(f'>0\) 可推出严格递增，但严格递增只能反推 \(f'\ge0\)，不能反推 \(f'>0\)。函数 \(x^3\) 在实线上严格递增而 \(f'(0)=0\)。
\end{remark}

\section{L'Hôpital 法则的各类严格版本}
\label{sec:lhopital-rules}

L'Hôpital 法则不是把“零除以零”或“无穷除以无穷”当作数值运算。它用 Cauchy 中值定理把函数增量比转成导数比，再由导数比的极限控制原商。为同时包括有限单侧极限和无穷远极限，先固定一种端点语言。

设 \(I\) 是非空开区间，它的两个端点允许属于
\(\overline{\R}=\R\cup\{-\infty,+\infty\}\)。以下写
\[
 x\to\alpha,\qquad x\in I,
\]
总表示 \(x\) 沿 \(I\) 趋向端点 \(\alpha\)。若 \(\alpha\) 有限，这是一个单侧极限；若 \(\alpha=\pm\infty\)，就是相应的无穷远极限。双侧有限极限应分别在左右两个区间上应用定理。

\subsection{导数比与割线比}

\begin{lemma}[Cauchy 比较引理]
\label{lem:cauchy-quotient-comparison}
设 \(f,g\) 在开区间 \(I\) 内可导，并且
\[
 g'(x)\ne0\qquad(x\in I).
\]
若 \(x,y\in I\)、\(x\ne y\)，则 \(g(x)\ne g(y)\)，且存在介于 \(x,y\) 之间的 \(\xi\)，使
\[
 \frac{f(y)-f(x)}{g(y)-g(x)}
 =\frac{f'(\xi)}{g'(\xi)}.
\]
因此，若导数比在连接 \(x,y\) 的区间内落在 \((A,B)\) 中，则相应增量比也落在 \((A,B)\) 中。
\end{lemma}

\begin{proof}
若 \(g(x)=g(y)\)，Rolle 定理给出介于两点之间的 \(\eta\) 使 \(g'(\eta)=0\)，与假设矛盾。把 Cauchy 中值定理用于以 \(x,y\) 为端点的闭区间，再除以非零的
\[
 \bigl(g(y)-g(x)\bigr)g'(\xi),
\]
便得商式及其比较结论。
\end{proof}

定理只需 \(g'\) 在所考察端点的某个邻域内不为零。实际应用时可以先缩短区间 \(I\)，所以下文为简化陈述写成 \(g'\ne0\) 于 \(I\)。

\subsection{\texorpdfstring{\(0/0\)}{0/0} 型}

\begin{theorem}[L'Hôpital 法则：\(0/0\) 型]
\label{thm:lhopital-zero-zero}
设 \(\alpha\) 是开区间 \(I\) 的一个端点，函数 \(f,g:I\to\R\) 在 \(I\) 内可导，并且
\[
 \lim_{\substack{x\to\alpha\\x\in I}}f(x)=0,
 \qquad
 \lim_{\substack{x\to\alpha\\x\in I}}g(x)=0,
 \qquad
 g'(x)\ne0\quad(x\in I).
\]
若导数比在扩充实数中有极限
\[
 \lim_{\substack{x\to\alpha\\x\in I}}
 \frac{f'(x)}{g'(x)}=L
 \qquad(L\in\overline{\R}),
\]
则 \(f(x)/g(x)\) 在充分靠近 \(\alpha\) 时有定义，并且
\[
 \lim_{\substack{x\to\alpha\\x\in I}}
 \frac{f(x)}{g(x)}=L.
\]
\end{theorem}

\begin{proof}
由比较引理，\(g\) 在 \(I\) 上为单射，因而至多有一个零点。若有这样的零点，它位于区间内部；把区间缩短到更靠近端点 \(\alpha\) 的一段即可避开它。因此 \(f/g\) 最终有定义。

先设 \(L\in\R\)。给定 \(\varepsilon>0\)，把导数比限制在端点的一个邻域内，使
\[
 L-\frac{\varepsilon}{2}
 <
 \frac{f'(t)}{g'(t)}
 <
 L+\frac{\varepsilon}{2}.
\]
固定这个邻域内的 \(x\)，并让 \(y\in I\) 位于 \(x\) 与端点 \(\alpha\) 之间。只要 \(y\) 也充分靠近 \(\alpha\)，连接 \(x,y\) 的整段区间都在上述邻域内。比较引理给出
\[
 L-\frac{\varepsilon}{2}
 <
 \frac{f(y)-f(x)}{g(y)-g(x)}
 <
 L+\frac{\varepsilon}{2}.
\]
令 \(y\to\alpha\)。因为 \(f(y)\to0\)、\(g(y)\to0\)，且已保证 \(g(x)\ne0\)，中间的增量比趋于
\[
 \frac{-f(x)}{-g(x)}=\frac{f(x)}{g(x)}.
\]
于是
\[
 L-\frac{\varepsilon}{2}
 \le\frac{f(x)}{g(x)}
 \le L+\frac{\varepsilon}{2},
\]
从而 \(\abs{f(x)/g(x)-L}<\varepsilon\)。这对所有充分靠近端点的 \(x\) 成立。

若 \(L=+\infty\)，给定 \(M>0\)，先使 \(f'/g'>M+1\)。同样的增量比估计在令 \(y\to\alpha\) 后给出
\[
 \frac{f(x)}{g(x)}\ge M+1>M.
\]
故商趋于 \(+\infty\)。若 \(L=-\infty\)，则 \((-f)'/g'\to+\infty\)。对 \(-f\) 使用刚证的情形，得到
\[
 \frac{-f(x)}{g(x)}\longrightarrow+\infty,
\]
故 \(f(x)/g(x)\to-\infty\)。
\end{proof}

上述证明在导数比尚无极限时仍给出定量夹逼。

\begin{corollary}[\(0/0\) 型的导数比夹逼]
\label{cor:lhopital-zero-zero-bracket}
在 \theoremref{thm:lhopital-zero-zero} 的条件中，用下列假设代替导数比极限：存在 \(m,M\in\R\) 和端点 \(\alpha\) 的一个邻域，使
\[
 m\le\frac{f'(t)}{g'(t)}\le M
\]
对该邻域内全部 \(t\in I\) 成立。则当 \(x\) 充分靠近 \(\alpha\) 时，
\[
 m\le\frac{f(x)}{g(x)}\le M.
\]
\end{corollary}

\begin{proof}
先缩短端点邻域，使其中 \(g\) 没有零点。固定邻域内的 \(x\)，并取位于 \(x\) 与端点之间的 \(y\)。Cauchy 比较引理给出
\[
 m\le\frac{f(y)-f(x)}{g(y)-g(x)}\le M.
\]
令 \(y\to\alpha\)，分子、分母分别趋于 \(-f(x)\)、\(-g(x)\)，中间商趋于 \(f(x)/g(x)\)。闭区间 \([m,M]\) 在取极限时保持，结论成立。
\end{proof}

令 \(m=L-\varepsilon\)、\(M=L+\varepsilon\)，就恢复有限极限版本。更一般地，这个夹逼说明即使导数比只在两个常数之间振荡，原商仍会受到相同的最终控制；但除非上下界可以任意收紧到同一个数，不能据此断言原商收敛。

有限端点的常见版本通常把 \(f,g\) 补定义到端点 \(\alpha\)，要求它们在那里连续且 \(f(\alpha)=g(\alpha)=0\)。这恰好保证上述两个单侧零极限；证明并不需要端点处可导。

\subsection{分母无界型}

\begin{theorem}[L'Hôpital 法则：无界分母型]
\label{thm:lhopital-infinite-denominator}
设 \(\alpha\) 是开区间 \(I\) 的一个端点，\(f,g:I\to\R\) 在 \(I\) 内可导，且
\[
 \abs{g(x)}\longrightarrow+\infty
 \qquad(x\to\alpha,\ x\in I),
\]
\[
 g'(x)\ne0\qquad(x\in I).
\]
若
\[
 \frac{f'(x)}{g'(x)}\longrightarrow L
 \qquad(x\to\alpha,\ x\in I)
\]
其中 \(L\in\overline{\R}\)，则
\[
 \frac{f(x)}{g(x)}\longrightarrow L.
\]
\end{theorem}

\begin{proof}
由 \(\abs{g(x)}\to+\infty\)，商在端点附近有定义。先设 \(L\in\R\)。给定 \(\varepsilon>0\)，把 \(c\in I\) 选入使
\[
 \abs{\frac{f'(t)}{g'(t)}-L}<\frac{\varepsilon}{4}
\]
对所有比 \(c\) 更靠近端点的 \(t\) 成立的邻域。当 \(x\) 比 \(c\) 更靠近端点时，比较引理给出
\[
 Q(x):=\frac{f(x)-f(c)}{g(x)-g(c)}
 =\frac{f'(\xi_x)}{g'(\xi_x)}
\]
其中 \(\xi_x\) 位于 \(c,x\) 之间。整段 \([c,x]\) 留在上述邻域，因此
\[
 \abs{Q(x)-L}<\frac{\varepsilon}{4}.
\]
恒等式
\[
 \frac{f(x)}{g(x)}
 =
 Q(x)\left(1-\frac{g(c)}{g(x)}\right)
 +\frac{f(c)}{g(x)}
\]
还可改写为
\[
 \frac{f(x)}{g(x)}-L
 =(Q(x)-L)\left(1-\frac{g(c)}{g(x)}\right)
 +\frac{f(c)-Lg(c)}{g(x)}.
\]
再把 \(x\) 限制得更靠近端点，使
\[
 \abs{\frac{g(c)}{g(x)}}<1,\qquad
 \abs{\frac{f(c)-Lg(c)}{g(x)}}<\frac{\varepsilon}{2}.
\]
于是右端绝对值小于
\[
 \frac{\varepsilon}{4}\cdot2+\frac{\varepsilon}{2}
 =\varepsilon,
\]
这就按量词证明了 \(f(x)/g(x)\to L\)。

若 \(L=+\infty\)，给定 \(M>0\)，把 \(c\) 取得充分靠近端点，使整段上的导数比大于 \(2(M+1)\)。于是更靠近端点时
\[
 Q(x)>2(M+1).
\]
再缩小端点邻域，使
\[
 \abs{\frac{g(c)}{g(x)}}<\frac12,
 \qquad
 \abs{\frac{f(c)}{g(x)}}<1.
\]
由上面的恒等式，
\[
 \frac{f(x)}{g(x)}
 >2(M+1)\cdot\frac12-1=M.
\]
故商趋于 \(+\infty\)。若 \(L=-\infty\)，对 \(-f\) 应用刚证的正无穷情形，得到
\[
 \frac{-f(x)}{g(x)}\longrightarrow+\infty,
\]
也就是 \(f(x)/g(x)\to-\infty\)。
\end{proof}

经典的“\(\infty/\infty\) 型”在此定理中还假设
\(\abs{f(x)}\to+\infty\)。证明实际上只使用了分母的无界趋势；分子的无界性不是结论所需的条件。因此“无界分母型”是更准确的定理名称。

\begin{corollary}[单侧、双侧与无穷远版本]
\label{cor:lhopital-directions}
上述两个定理适用于有限点的左极限、右极限，以及 \(x\to+\infty\) 或 \(x\to-\infty\)。若要推出有限点的双侧极限，应分别在左右两侧核对假设；只有两侧所得极限相同，双侧结论才成立。导数比极限可以是有限实数，也可以是 \(+\infty\) 或 \(-\infty\)。
\end{corollary}

\begin{example}[\(0/0\) 与单侧极限]
假定正弦函数的导数公式已经由独立方法建立。在 \(x\to0\) 时，
\[
 \frac{\sin x}{x}
\]
为 \(0/0\) 型。分母导数恒为 \(1\)，导数比为 \(\cos x\to1\)，故商趋于 \(1\)。这个双侧结论可以看作分别应用于 \((-\delta,0)\) 与 \((0,\delta)\) 的结果。
\end{example}

若所采用的正弦函数求导证明本身以
\(\lim_{x\to0}\sin x/x=1\) 为前提，就不能再用本例反向证明该极限；这会形成循环论证。此时应把本例理解为 L'Hôpital 法则对一个已知极限的复核，或者改用不依赖该极限的方式先建立三角函数导数。

\begin{example}[无穷远与无界分母]
当 \(x\to+\infty\) 时，
\[
 \frac{\log x}{x}
\]
的分母趋于 \(+\infty\)，而
\[
 \frac{(\log x)'}{x'}=\frac1x\longrightarrow0.
\]
所以 \(\log x/x\to0\)。分子虽然也趋于 \(+\infty\)，证明只需分母无界。
\end{example}

\begin{example}[有限端点的 \(\infty/\infty\) 型]
当 \(x\to0+\) 时，\(\log x\to-\infty\)、\(1/x\to+\infty\)，且
\[
 \frac{(\log x)'}{(1/x)'}
 =\frac{1/x}{-1/x^2}=-x\longrightarrow0.
\]
因此
\[
 x\log x=\frac{\log x}{1/x}\longrightarrow0.
\]
\end{example}

\section{L'Hôpital 法则的条件与反例}
\label{sec:lhopital-counterexamples}

\subsection{定理只给出充分条件}

L'Hôpital 法则的逻辑形式是：在指定不定型和正则性条件下，导数比极限存在则函数比具有同一极限。它不是等价判据。函数比的极限可能存在，而导数比没有极限。

\begin{counterexample}[逆命题失败]
令 \(g(x)=x\)，并在 \(x\ne0\) 时令
\[
 f(x)=x+x^2\sin\frac1{x^2},
\qquad f(0)=0.
\]
当 \(x\to0\) 时，
\[
 \frac{f(x)}{g(x)}
 =1+x\sin\frac1{x^2}\longrightarrow1.
\]
但对 \(x\ne0\)，
\[
 \frac{f'(x)}{g'(x)}
 =1+2x\sin\frac1{x^2}
   -\frac2x\cos\frac1{x^2},
\]
它无界并且没有极限。因此函数比收敛不能反推导数比收敛。
\end{counterexample}

\begin{counterexample}[缺少不定型条件]
当 \(x\to0\) 时取
\[
 f(x)=1+x,\qquad g(x)=x.
\]
导数比恒为 \(1\)，但
\[
 \frac{f(x)}{g(x)}=\frac1x+1
\]
没有有限双侧极限，更不趋于 \(1\)。这里分子不趋于零，分母也不无界，所以两个 L'Hôpital 定理都不适用。
\end{counterexample}

\subsection{导数比必须有确定极限}

导数比有界不足以推出函数比收敛。令 \(x>1\)，
\[
 g(x)=x,\qquad f(x)=x\bigl(2+\sin(\log x)\bigr).
\]
则 \(f(x),g(x)\to+\infty\)，而
\[
 \frac{f(x)}{g(x)}=2+\sin(\log x)
\]
没有极限。相应导数比
\[
 \frac{f'(x)}{g'(x)}
 =2+\sin(\log x)+\cos(\log x)
\]
虽然有界，也没有极限。L'Hôpital 法则要求导数比本身趋于某个扩充实数，不能用“有界”或“沿某条数列收敛”替代。

条件 \(g'\ne0\) 的作用是保证所有相关割线分母非零，并使导数比在完整端点邻域内有定义。它可以在更精细的版本中削弱，但必须用别的条件保证 Cauchy 增量比论证有效。端点本身的 \(g'(\alpha)\) 不参与定理；例如 \(g(x)=x^2\) 在 \(x\to0\) 时虽然 \(g'(0)=0\)，但在每个删心侧邻域内 \(g'(x)\ne0\)，仍可使用定理。

\subsection{单侧核对与重复使用}

\begin{example}[两侧结果不同]
函数
\[
 \frac{\abs x}{x}
\]
在 \(x\to0+\) 时的导数比为 \(1\)，在 \(x\to0-\) 时的导数比为 \(-1\)。分子在 \(0\) 处不可导不影响两个删心侧区间内的应用；两侧结论不同，所以双侧极限不存在。把一个单侧证明直接写成双侧证明会遗漏这一差异。
\end{example}

重复使用 L'Hôpital 法则时，每一次都要重新检查所得商仍属于定理覆盖的情形。例如
\[
 \lim_{x\to0}\frac{1-\cos x}{x^2}
\]
第一次化为
\[
 \lim_{x\to0}\frac{\sin x}{2x},
\]
它仍是 \(0/0\) 型；第二次化为
\[
 \lim_{x\to0}\frac{\cos x}{2}=\frac12.
\]
若中间商已经不再是 \(0/0\) 型或无界分母型，不能只因为表达式仍“看起来复杂”就继续求导。

乘积 \(0\cdot\infty\)、差式 \(\infty-\infty\) 和幂式 \(1^\infty\) 不在定理的直接陈述中。可以先把它们改写为商或取对数，但改写以后仍须核对定义域、符号以及两个基本定理的全部假设。代数改写不会自动产生可用的 L'Hôpital 条件。

\section{向量值映射中值定理的局限}
\label{sec:vector-mean-value}

本节在 \(\R^m\) 上使用标准欧氏内积 \(\langle\cdot,\cdot\rangle\) 及其诱导范数 \(\norm{\cdot}\)。

标量证明能够从局部极值推出导数为零，并能用 Rolle 定理找到一个同时满足全部要求的中间点。向量值函数没有与实数次序相容的最大值、最小值；各坐标使用标量中值定理时，得到的中间点也可能不同。

\begin{counterexample}[不存在一般的向量等式]
\label{ex:vector-mvt-failure}
令
\[
 F(t)=(\cos t,\sin t),
\qquad 0\le t\le2\pi.
\]
有 \(F(2\pi)-F(0)=0\)。若存在标量中值定理形式的点 \(c\in(0,2\pi)\) 使
\[
 F(2\pi)-F(0)=F'(c)\,2\pi,
\]
就会有 \(F'(c)=0\)。但
\[
 F'(t)=(-\sin t,\cos t),\qquad \norm{F'(t)}=1,
\]
矛盾。因此即使 \(F\) 光滑，也未必存在一个共同中间点再现整个向量增量。
\end{counterexample}

逐坐标使用 Lagrange 定理只能得到
\[
 F_j(b)-F_j(a)=F_j'(c_j)(b-a),
\]
其中 \(c_j\) 可能随坐标 \(j\) 改变。这组等式不能合并成一个共同的 \(c\)。

\begin{theorem}[向量值中值估计]
\label{thm:vector-mean-value-estimate}
设 \(a<b\)，\(F:[a,b]\to\R^m\) 连续并在 \((a,b)\) 内可导，其中导数按各坐标定义为
\[
 F'(t)=\bigl(F_1'(t),\ldots,F_m'(t)\bigr).
\]
若
\[
 \norm{F'(t)}\le M\qquad(t\in(a,b)),
\]
则
\[
 \norm{F(b)-F(a)}\le M(b-a).
\]
更精确地，若 \(F(b)\ne F(a)\)，则存在 \(c\in(a,b)\) 使
\[
 \norm{F(b)-F(a)}
 \le (b-a)\norm{F'(c)}.
\]
\end{theorem}

\begin{proof}
若 \(F(b)=F(a)\)，结论显然。否则令
\[
 v=\frac{F(b)-F(a)}{\norm{F(b)-F(a)}},
\]
并考虑标量函数
\[
 \varphi(t)=\langle v,F(t)\rangle.
\]
它在 \([a,b]\) 上连续、在内部可导，且
\[
 \varphi'(t)=\langle v,F'(t)\rangle.
\]
Lagrange 中值定理给出某个 \(c\in(a,b)\)，使
\begin{align*}
 \norm{F(b)-F(a)}
 &=\langle v,F(b)-F(a)\rangle\\
 &=(b-a)\langle v,F'(c)\rangle\\
 &\le(b-a)\norm{F'(c)}.
\end{align*}
最后使用 \(\norm{F'(c)}\le M\) 即得统一估计。
\end{proof}

\begin{corollary}[向量值 Lipschitz 判据]
若 \(F:I\to\R^m\) 在区间 \(I\) 的内点可导，并满足
\[
 \norm{F'(t)}\le M,
\]
且在属于 \(I\) 的端点连续，则
\[
 \norm{F(y)-F(x)}\le M\abs{y-x}
\qquad(x,y\in I).
\]
\end{corollary}

\begin{proof}
对以 \(x,y\) 为端点的闭区间应用向量值中值估计。
\end{proof}

因此，从标量中值定理推广到向量值映射时，应保留的是增量的范数估计，而不是强求一个中间点给出向量等式。证明所用的方向 \(v\) 由端点增量决定；对其他投影方向，标量中值定理可能给出不同的中间点。

\section*{习题}
\addcontentsline{toc}{section}{习题}

\subsection*{A组　中值定理与基本推论}
\begin{enumerate}[label=\textbf{\thechapter.\arabic*.},ref=\thechapter.\arabic*,leftmargin=3.4em,labelsep=0.7em,itemsep=0.85em,topsep=0.8em]
 \exerciseitem{fermat-local-extremum} 从左右差商的符号出发完整证明 Fermat 定理，并说明为什么证明不能用于区间端点。
 \exerciseitem{rolle-polynomial-root} 设实多项式 \(p\) 有两个不同实根。证明 \(p'\) 在两根之间至少有一个实根；进而证明 \(p\) 有 \(k\) 个不同实根时，\(p'\) 至少有 \(k-1\) 个不同实根。
 \exerciseitem{lagrange-secant-example} 对 \(f(x)=x^3\) 在区间 \([a,b]\) 上写出 Lagrange 中值点所满足的方程，并在 \([0,1]\) 上求出全部中值点。
 \exerciseitem{zero-derivative-constant} 证明区间上处处可导且导数恒为零的函数为常值，并构造不连通定义域上的反例。
 \exerciseitem{derivative-sign-monotonicity} 证明若 \(f'\ge0\) 于区间内部，则 \(f\) 单调不减；若 \(f'>0\)，则严格递增。
 \exerciseitem{nonnegative-derivative-strictness} 设 \(f'\ge0\) 于区间内部。证明 \(f\) 严格递增当且仅当不存在非退化子区间使 \(f'\) 在其内部恒为零；给出一个导数在无穷多个点为零但函数仍严格递增的例子。
 \exerciseitem{monotone-derivative-converse} 证明可导的单调不减函数满足 \(f'\ge0\)，并举例说明严格递增函数的导数仍可在某点为零。
 \exerciseitem{derivative-bounds-increments} 设 \(m\le f'\le M\)。证明对 \(x<y\) 有
 \[
  m(y-x)\le f(y)-f(x)\le M(y-x).
 \]
 \exerciseitem{bounded-derivative-lipschitz} 证明 \(\abs{f'}\le K\) 推出 \(f\) 为 \(K\)-Lipschitz，并说明闭区间端点处为何不需要导数。
 \exerciseitem{lipschitz-derivative-bound} 设 \(f\) 在区间上为 \(K\)-Lipschitz 且在内点 \(a\) 可导。证明 \(\abs{f'(a)}\le K\)。
 \exerciseitem{cauchy-mvt-ratio} 从交叉相乘形式推出 Cauchy 中值定理的增量比形式，逐项说明两个非零分母由何处得到。
 \exerciseitem{derivative-comparison} 若 \(f'\le g'\) 于区间内部，证明
 \[
  f(y)-f(x)\le g(y)-g(x)\qquad(x<y).
 \]
 \exerciseitem{endpoint-derivative-unneeded} 以 \(\sqrt x\) 在 \([0,1]\) 上的限制为例，核对 Lagrange 中值定理全部假设并求一个中值点，说明端点可导不是必要条件。
\end{enumerate}

\subsection*{B组　条件辨析与 L'Hôpital 法则}
\begin{enumerate}[label=\textbf{\thechapter.\arabic*.},ref=\thechapter.\arabic*,leftmargin=3.4em,labelsep=0.7em,itemsep=0.85em,topsep=0.8em,resume]
 \exerciseitem{fermat-hypotheses-counterexamples} 分别构造反例说明 Fermat 定理中的内点、可导性和局部极值三个条件不能删去或反向使用。
 \exerciseitem{rolle-hypotheses-counterexamples} 分别删去 Rolle 定理的闭区间连续性、内部可导性和端点等值条件，给出结论失败的反例，并核对其余假设。
 \exerciseitem{interval-hypothesis-counterexample} 构造不连通定义域上的函数，使两个指定端点函数值相同而所有可导点的导数均不为零；解释区间性在 Rolle 证明中的位置。
 \exerciseitem{strict-versus-nonstrict-monotonicity} 分别给出下列反向命题的反例：单调不减推出 \(f'>0\)；严格递增推出 \(f'>0\)；\(f'\ge0\) 推出严格递增。
 \exerciseitem{inverse-lipschitz-from-lower-derivative} 设 \(f'\ge m>0\) 于 \((a,b)\)。证明 \(f^{-1}\) 在像区间上为 \(1/m\)-Lipschitz；指出只假设 \(f'>0\) 而没有统一下界时为何得不到固定常数。
 \exerciseitem{lhopital-zero-zero-sine} 假定正弦函数的导数公式已由不依赖下述极限的方法建立。用 \(0/0\) 型 L'Hôpital 法则证明
 \[
  \lim_{x\to0}\frac{\sin x}{x}=1,
 \]
 并把双侧证明拆成两个单侧应用。
 \exerciseitem{lhopital-zero-zero-infinity} 用无穷远处的 \(0/0\) 型定理计算
 \[
  \lim_{x\to+\infty}\frac{1/x}{1/\sqrt{x}}.
 \]
 \exerciseitem{lhopital-infinity-infinity-log} 用无界分母型定理证明
 \[
  \lim_{x\to+\infty}\frac{\log x}{x}=0.
 \]
 \exerciseitem{lhopital-infinite-result} 用 L'Hôpital 法则证明
 \[
  \lim_{x\to+\infty}\frac{e^x-1}{x}=+\infty,
 \]
 并核对无穷导数比版本中的量词。
 \exerciseitem{lhopital-nonindeterminate-failure} 构造导数比恒有有限极限，但原商因既非 \(0/0\) 型也非无界分母型而不趋于该极限的例子。
 \exerciseitem{lhopital-converse-failure} 验证正文中的函数
 \[
  f(x)=x+x^2\sin(1/x^2),\qquad g(x)=x
 \]
 满足 \(f/g\to1\) 而 \(f'/g'\) 没有极限。
 \exerciseitem{lhopital-repeated-application} 假定三角函数导数公式已经独立建立。严格核对两次使用 L'Hôpital 法则计算
 \[
  \lim_{x\to0}\frac{1-\cos x}{x^2}
 \]
 时每一步的假设。
 \exerciseitem{lhopital-derivative-ratio-bracket} 在 \(0/0\) 型 L'Hôpital 法则的条件下，设端点附近
 \[
  m\le\frac{f'}{g'}\le M.
 \]
 不假设导数比有极限，直接由 Cauchy 比较引理证明原商最终也位于 \([m,M]\)。再用这一结论证明：若对每个 \(\varepsilon>0\)，端点附近都有
 \[
  L-\varepsilon\le\frac{f'}{g'}\le L+\varepsilon,
 \]
 则 \(f/g\to L\)。
\end{enumerate}

\subsection*{C组　比较、无穷远与向量值映射}
\begin{enumerate}[label=\textbf{\thechapter.\arabic*.},ref=\thechapter.\arabic*,leftmargin=3.4em,labelsep=0.7em,itemsep=0.85em,topsep=0.8em,resume]
 \exerciseitem{cauchy-comparison-full-proof} 设 \(g'\ne0\)，并有
 \[
  m\le f'/g'\le M.
 \]
 不跳步地证明任意两点之间的增量比也位于 \([m,M]\)，包括 \(g(y)-g(x)\ne0\) 的证明。
 \exerciseitem{lhopital-left-and-two-sided} 写出 \(x\to a-\) 的 \(0/0\) 型 L'Hôpital 定理；再证明若左右导数比极限均为同一 \(L\in\overline{\R}\)，则原商有双侧极限 \(L\)。
 \exerciseitem{lhopital-at-negative-infinity} 把无界分母型定理完整改写为 \(x\to-\infty\) 的版本，并用它计算
 \[
  \lim_{x\to-\infty}\frac{\log(1+x^2)}{x}.
 \]
 \exerciseitem{lhopital-denominator-divergence-suffices} 给出一个分子具有有限极限、分母绝对值趋于无穷的例子，用无界分母型定理求商极限，并说明为何该定理不必假设分子无界。
 \exerciseitem{lhopital-bounded-derivative-ratio-insufficient} 验证
 \[
  f(x)=x(2+\sin\log x),\qquad g(x)=x
 \]
 在 \(x\to+\infty\) 时构成 \(\infty/\infty\) 型，导数比有界而函数比没有极限。
 \exerciseitem{vector-mvt-circle-counterexample} 对 \(F(t)=(\cos t,\sin t)\) 证明不存在
 \[
  F(2\pi)-F(0)=2\pi F'(c)
 \]
 的中间点 \(c\)。
 \exerciseitem{vector-mvt-norm-estimate} 通过端点增量方向的标量投影，完整证明向量值中值估计。
 \exerciseitem{vector-derivative-bound-lipschitz} 设 \(F:I\to\R^m\) 可导且 \(\norm{F'}\le M\)。证明 \(F\) 为 \(M\)-Lipschitz，并说明这个常数在 \(F(t)=tv\) 时达到。
 \exerciseitem{vector-projection-mean-value} 设 \(a<b\)，\(F:[a,b]\to\R^m\) 连续并在 \((a,b)\) 内可导。证明对每个固定 \(u\in\R^m\)，存在依赖于 \(u\) 的 \(c_u\in(a,b)\)，使
 \[
  \langle u,F(b)-F(a)\rangle
  =(b-a)\langle u,F'(c_u)\rangle.
 \]
 解释为什么这些 \(c_u\) 一般不能取成同一点。
 \projectexerciseitem{mean-value-theorem-dependency-project} 【前置：本章；预计用时：90分钟】整理
 \[
  \text{极值定理与 Fermat 定理}\Longrightarrow\text{Rolle}
  \Longrightarrow\text{Lagrange}
  \Longrightarrow\text{Cauchy}
  \Longrightarrow\text{L'Hôpital}
 \]
 的证明依赖。对每个箭头写出辅助函数或关键估计，并给出至少一个删条件反例；最后说明向量值情形在哪个箭头失去标量等式。
\end{enumerate}
